\documentclass[aspectratio=169,11pt]{beamer}

\usetheme{Madrid}
\usecolortheme{default}

\usepackage[utf8]{inputenc}
\usepackage[T1]{fontenc}
\usepackage{booktabs}
\usepackage{array}

\definecolor{garpblue}{RGB}{0,51,102}
\definecolor{garpgold}{RGB}{204,153,0}

\setbeamercolor{structure}{fg=garpblue}
\setbeamercolor{title}{fg=white,bg=garpblue}
\setbeamercolor{frametitle}{fg=white,bg=garpblue}
\setbeamercolor{block title}{fg=white,bg=garpblue}
\setbeamercolor{block body}{bg=garpblue!5}
\setbeamercolor{itemize item}{fg=garpgold}
\setbeamercolor{itemize subitem}{fg=garpblue}

\hypersetup{
	colorlinks=true,
	linkcolor=garpblue,
	urlcolor=garpgold,
	pdftitle={GARP Risk and AI Exam Prep Course - Unofficial},
	pdfauthor={Unofficial Course Notes}
}

\title[GARP Risk \& AI]{GARP Risk and AI Exam Prep Course}
\subtitle{Unofficial Exam Preparation for AI in Risk Management}
\author[Unofficial]{Unofficial Course Notes}
\institute[GARP Risk \& AI]{6 hours of video content \\ + PDF notes, question banks, and tutorials}
\date{\today}

\begin{document}
	
	\begin{frame}
		\titlepage
	\end{frame}
	
	\begin{frame}{Outline}
		\tableofcontents
	\end{frame}
	
	\section{Intended Learners}
	
	\begin{frame}{Intended Learners --- Course Requirements}
		\begin{block}{Prerequisites and Requirements}
			\begin{itemize}
				\item Basic understanding of risk management concepts (helpful but not mandatory)
				\item Familiarity with statistics and probability at an introductory level
				\item Basic knowledge of machine learning terminology
				\item No programming experience required
				\item A willingness to work through PDF notes, question banks, and tutorials
			\end{itemize}
		\end{block}
	\end{frame}
	
	\begin{frame}{Intended Learners --- Learning Objectives}
		\begin{block}{By the end of this course, you will be able to:}
			\begin{enumerate}
				\item Explain how artificial intelligence is applied within risk management frameworks.
				\item Identify and compare key AI and ML tools and techniques.
				\item Evaluate model performance and estimation methods.
				\item Describe responsible and ethical AI principles and governance.
				\item Apply generative AI and NLP concepts to risk use cases.
				\item Answer exam-style questions across all five modules confidently.
			\end{enumerate}
		\end{block}
	\end{frame}
	
	\begin{frame}{Intended Learners --- Who This Course Is For}
		\begin{block}{Primary Audience}
			\begin{itemize}
				\item Risk analysts and risk managers
				\item Data scientists entering risk domains
				\item AI and ML practitioners in financial services
				\item Compliance and governance professionals
			\end{itemize}
		\end{block}
		\begin{block}{Secondary Audience}
			\begin{itemize}
				\item Students preparing for GARP Risk and AI certification
				\item Auditors reviewing AI models
				\item Consultants advising on AI risk
				\item Anyone curious about AI in risk management
			\end{itemize}
		\end{block}
	\end{frame}
	
	\section{Curriculum}
	
	\begin{frame}{Curriculum Overview}
		\begin{block}{Course Structure}
			\begin{itemize}
				\item 6 Sections with 30-40 Lectures
				\item 6-10 hrs of video content
				\item Each module includes notes, questions, and tutorials in PDF format
				\item Aligned with the official GARP Risk and AI curriculum
			\end{itemize}
		\end{block}
		\begin{block}{Note}
			All lectures have accompanying content. No lecture is left without supporting materials.
		\end{block}
	\end{frame}
	
	\begin{frame}{Section 1 and 2: Introduction and Module 1}
		\begin{block}{Section 1: Introduction}
			\begin{itemize}
				\item Lecture 1: Introduction
			\end{itemize}
		\end{block}
		\begin{block}{Section 2: Module 1 --- AI and Risk: Introduction and Overview}
			\begin{itemize}
				\item Lecture 2: Module 1 AI and Risk --- Introduction and Overview (Preview enabled)
				\item Lecture 3: Questions --- Module 1 AI and Risk --- Introduction and Overview
			\end{itemize}
		\end{block}
	\end{frame}
	
	\begin{frame}{Section 3: Module 2 --- Tools and Techniques (Part 1)}
		\begin{block}{Lectures 4 to 10}
			\begin{itemize}
				\item Lecture 4: Chapter 1 --- Introduction to Tools and Techniques (Preview enabled)
				\item Lecture 5: Chapter 2 --- Unsupervised Learning
				\item Lecture 6: Chapter 3 --- Supervised Learning for Numerical Data
				\item Lecture 7: Chapter 4 --- Decision Trees
				\item Lecture 8: Chapter 4 --- KNN Theory
				\item Lecture 9: Chapter 5 --- Semi-Supervised Learning
				\item Lecture 10: Chapter 6 --- Reinforcement Learning
			\end{itemize}
		\end{block}
	\end{frame}
	
	\begin{frame}{Section 3: Module 2 --- Tools and Techniques (Part 2)}
		\begin{block}{Lectures 11 to 16}
			\begin{itemize}
				\item Lecture 11: Chapter 7 --- Supervised Learning: Model Estimation
				\item Lecture 12: Chapter 8 --- Supervised Learning: Model Performance
				\item Lecture 13: Chapter 9 --- Natural Language Processing Part 1
				\item Lecture 14: Chapter 10 --- Generative AI Part 1
				\item Lecture 15: Chapter 10 --- Generative AI Part 2
				\item Lecture 16: Chapter 10 --- Generative AI Part 3
			\end{itemize}
		\end{block}
	\end{frame}
	
	\begin{frame}{Section 4, 5 and 6: Modules 3 to 5}
		\begin{block}{Section 4: Module 3 --- Risk and Risk Factors}
			\begin{itemize}
				\item Lecture 17: Module 3 --- Risk and Risk Factors
				\item Lecture 18: Questions --- Risk and Risk Factors
			\end{itemize}
		\end{block}
		\begin{block}{Section 5: Module 4 --- Responsible and Ethical AI}
			\begin{itemize}
				\item Lecture 19: Module 4 --- Responsible and Ethical AI
			\end{itemize}
		\end{block}
		\begin{block}{Section 6: Module 5 --- Data and AI Model Governance}
			\begin{itemize}
				\item Lecture 20: Module 5 --- Data and AI Model Governance
				\item Lecture 21: Questions --- Data and AI Model Governance
			\end{itemize}
		\end{block}
	\end{frame}
	
	\section{Course Landing Page}
	
	\begin{frame}{Course Landing Page --- Description (200+ words)}
		\begin{block}{Course Description}
			\small
			This unofficial GARP Risk and AI Exam Prep Course is designed for professionals and students who want to master the intersection of artificial intelligence and risk management. Whether you are a risk analyst, data scientist, compliance officer, or AI practitioner, this course provides a structured path through the key concepts, tools, and governance frameworks required by the GARP Risk and AI certification.
			
			The course includes 4 hours and 20 minutes of video content, organized into six sections and twenty-one lectures. You will explore AI and risk fundamentals, supervised and unsupervised learning, decision trees, KNN, semi-supervised learning, reinforcement learning, model estimation, model performance evaluation, natural language processing, and generative AI. Each module is supported by comprehensive PDF notes, question banks, tutorials, and variation materials to reinforce your learning.
			
			Beyond technical content, the course covers responsible and ethical AI, data governance, and AI model governance. These are critical topics for anyone deploying AI in regulated environments. You will learn how to evaluate model performance, identify risks, and apply governance frameworks to real-world scenarios.
			
			The course is self-paced, beginner-friendly, and requires no programming experience. All lectures have accompanying content, and the included question PDFs help you test your understanding after each module. By the end, you will be prepared to answer exam-style questions confidently and apply AI risk concepts in your daily work.
			
			This is an unofficial study aid and is not affiliated with GARP or Udemy. It is intended for personal exam preparation and professional development.
		\end{block}
	\end{frame}
	
	\begin{frame}{Course Landing Page --- Subtitle, Category, Level}
		\begin{block}{Course Subtitle}
			Unofficial Exam Preparation for AI in Risk Management
		\end{block}
		\begin{block}{Category}
			Business --- Business Analytics and Intelligence
		\end{block}
		\begin{block}{Level}
			Intermediate
		\end{block}
		\begin{block}{Subcategory}
			Risk Management
		\end{block}
		\begin{block}{Primarily Taught}
			Artificial Intelligence (AI) and Risk Management
		\end{block}
	\end{frame}
	
	\begin{frame}{Course Landing Page --- Course Image}
		\begin{block}{Course Image Requirements}
			\begin{itemize}
				\item Recommended size: 750 by 422 pixels (16:9 aspect ratio)
				\item Format: JPG or PNG
				\item Content: AI and risk visual (neural network, shield, risk charts)
				\item Text: Avoid too much text; keep it clean and professional
				\item Branding: Use dark blue and gold to match the course theme
			\end{itemize}
		\end{block}
		\begin{block}{Note}
			Upload the course image on the Course Landing Page before publishing.
		\end{block}
	\end{frame}
	
	\section{Available PDF Materials}
	
	\begin{frame}{Chapter PDFs (Part 1)}
		\small
		\centering
		\begin{tabular}{@{}p{0.50\textwidth}p{0.42\textwidth}@{}}
			\toprule
			File & Description \\
			\midrule
			Chapter1Introduction.pdf & Introduction to AI and risk \\
			Chapter2UnsupervisedLearning.pdf & Unsupervised learning \\
			Chapter3supervisedLearning\_v3.pdf & Supervised learning \\
			Chapter4Techniques.pdf & ML techniques (Part 2) \\
			Chapter4Techniques\_notes.pdf & Notes for Chapter 4 \\
			Chapter4Techniques\_variation.pdf & Variation of Chapter 4 \\
			Chapter5SemiSupervisedLearning.pdf & Semi-supervised learning \\
			Chapter5SemiSupervised\_notes.pdf & Notes for Chapter 5 \\
			Chapter5SemiSupervised\_v2.pdf & Version 2 of Chapter 5 \\
			Chapter6ReinforcementLearning.pdf & Reinforcement learning \\
			Chapter6RL\_tutorial.pdf & RL tutorial \\
			Chapter6RL\_tutorial\_v2.pdf & Tutorial version 2 \\
			\bottomrule
		\end{tabular}
	\end{frame}
	
	\begin{frame}{Chapter PDFs (Part 2)}
		\small
		\centering
		\begin{tabular}{@{}p{0.50\textwidth}p{0.42\textwidth}@{}}
			\toprule
			File & Description \\
			\midrule
			Chapter6RL\_tutorial\_v2\_part1.pdf & Tutorial v2, Part 1 \\
			Chapter6RL\_tutorial\_part1.pdf & Tutorial Part 1 (2026-09-05) \\
			Chapter6RL\_tutorial\_part2.pdf & Tutorial Part 2 (2026-09-05) \\
			Chapter7ModelEstimation.pdf & Model estimation \\
			Chapter7ModelEstimation\_notes.pdf & Notes for Chapter 7 \\
			Chapter7tutorial.pdf & Chapter 7 tutorial \\
			Chapter8PerformanceEvaluation.pdf & Model performance evaluation \\
			Chapter8Performance\_notes.pdf & Notes for Chapter 8 \\
			Chapter8Performance\_variation.pdf & Variation of Chapter 8 \\
			Chapter9NLP.pdf & Natural language processing \\
			Chapter9NLP\_tutorial.pdf & NLP tutorial \\
			\bottomrule
		\end{tabular}
	\end{frame}
	
	\begin{frame}{Module PDFs}
		\small
		\centering
		\begin{tabular}{@{}p{0.50\textwidth}p{0.42\textwidth}@{}}
			\toprule
			File & Description \\
			\midrule
			Chapter10Generative.pdf & Generative AI \\
			Chapter10Generative\_questions.pdf & Generative AI questions \\
			Module1\_notes.pdf & Module 1 notes \\
			Module1\_questions.pdf & Module 1 questions \\
			Module3RiskFactors.pdf & Risk and risk factors \\
			Module3RiskFactors\_notes.pdf & Module 3 notes \\
			Module3RiskFactors\_q.pdf & Module 3 questions \\
			Module4ResponsibleEthicalAI.pdf & Responsible and ethical AI \\
			Module4Responsible\_notes.pdf & Module 4 notes \\
			Module4Responsible\_q.pdf & Module 4 questions \\
			Module5Governance.pdf & Data and AI governance \\
			Module5Governance\_notes.pdf & Module 5 notes \\
			Module5Governance\_q.pdf & Module 5 questions \\
			\bottomrule
		\end{tabular}
	\end{frame}
	
	\section{How to Use This Course}
	
	\begin{frame}{How to Use This Course}
		\begin{enumerate}
			\item Watch the video lectures in the order listed under each section.
			\item Read the corresponding PDF notes for each chapter or module.
			\item Attempt the question PDFs after finishing each module.
			\item Work through the tutorials for Chapters 6, 7, and 9.
			\item Review the variation PDFs for alternative explanations.
			\item Use the 100-slide summary PDF for quick revision before the exam.
		\end{enumerate}
		\begin{block}{Note}
			The course also includes explanations with question-and-answer sections for selected modules.
		\end{block}
	\end{frame}
	
	\section{Disclaimer}
	
	\begin{frame}{Disclaimer}
		\begin{block}{Important}
			\begin{itemize}
				\item This document is an unofficial study aid.
				\item Not affiliated with, endorsed by, or sponsored by GARP or Udemy.
				\item All trademarks and copyrights belong to their respective owners.
				\item PDF filenames and course structure are reproduced for organizational and study purposes only.
				\item Always refer to the official GARP Risk and AI curriculum for the authoritative syllabus.
			\end{itemize}
		\end{block}
	\end{frame}
	
	\begin{frame}
		\centering
		\Huge \textcolor{garpblue}{Thank You!}
		\vspace{1cm}
		
		\Large Good luck with your GARP Risk and AI exam!
		\vspace{1cm}
		
		\normalsize
		\textit{Unofficial study aid --- not affiliated with GARP or Udemy.}
	\end{frame}
	
\end{document}